% !TeX root = ../深度学习与大语言模型技术原理.tex
% 由主文档通过 \input 载入；请编译主文档。
\setcounter{chapter}{2} % 保留原书章号。
\chapter{神经网络、自动微分与优化}
\label{ch:optimization}

\section{多层表示与反向传播}

多层感知机（multilayer perceptron, MLP）把线性映射与非线性激活复合：$h^{(\ell)}=\varphi(W^{(\ell)}h^{(\ell-1)}+b^{(\ell)})$。若每一层都是线性的，所有层仍可合并成单个线性映射；非线性使模型能够表达输入相关的分区和交互。常见激活包括 $\operatorname{ReLU}(x)=\max(0,x)$、$\operatorname{SiLU}(x)=x\sigma(x)$ 与 $\operatorname{GELU}(x)=x\Phi(x)$，其中 $\sigma(x)=(1+e^{-x})^{-1}$，$\Phi$ 为标准正态分布函数。

\emph{反向传播（backpropagation）}是计算图上的链式法则。对 $z=Wx+b$ 和标量损失 $\mathcal L$，若上游梯度为 $\delta=\partial\mathcal L/\partial z$，则
\[
\frac{\partial\mathcal L}{\partial W}=\delta x^\top,\qquad
\frac{\partial\mathcal L}{\partial b}=\delta,\qquad
\frac{\partial\mathcal L}{\partial x}=W^\top\delta.
\]
沿图逆序执行这些向量--雅可比积，可避免显式构造整个巨大雅可比矩阵。自动微分计算的是程序所表达函数的导数，无法自动修正标签移位、错误掩码或错误损失归一化。

深层网络的梯度含多个雅可比矩阵的乘积。其范数满足 $\|J_L\cdots J_1\|\le\prod_\ell\|J_\ell\|$；连续放大或衰减可能造成梯度爆炸或消失，但这个上界不意味着所有方向都有相同行为。残差连接 $h^{(\ell+1)}=h^{(\ell)}+F_\ell(h^{(\ell)})$ 使局部雅可比变为 $I+J_{F_\ell}$，提供直接传播路径；初始化、归一化和学习率仍然重要。

\section{随机梯度下降、动量与自适应更新}

设总体风险为 $\mathcal L(\theta)=\mathbb E_Z[\ell(\theta;Z)]$。在可以交换梯度与期望且批量按目标分布抽样时，批梯度 $g_t$ 是总体梯度的无偏估计。\emph{随机梯度下降（stochastic gradient descent, SGD）}用 $\theta_{t+1}=\theta_t-\eta_tg_t$ 更新。动量法以 $v_t=\mu v_{t-1}+g_t$ 累积方向，再用 $-\eta_tv_t$ 更新，从而平滑高频噪声，但不同软件可能采用带 $(1-\mu)$ 的另一种约定。

Adam 维护梯度的一阶和二阶矩估计，初始化 $m_0=v_0=0$：
\begin{align}
m_t&=\beta_1m_{t-1}+(1-\beta_1)g_t, &
v_t&=\beta_2v_{t-1}+(1-\beta_2)g_t^{\odot2},\\
\widehat m_t&=m_t/(1-\beta_1^t), &
\widehat v_t&=v_t/(1-\beta_2^t),\\
\theta_{t+1}&=\theta_t-\eta_t\frac{\widehat m_t}{\sqrt{\widehat v_t}+\varepsilon}.&&
\end{align}
平方、平方根和除法逐元素计算。偏差修正补偿零初始化导致的早期缩小；分母按坐标调节更新尺度，并不等于计算了 Hessian，也不保证非凸问题收敛到全局最优解。

\emph{AdamW} 把权重衰减从自适应梯度中解耦\cite{adamw}：
\begin{equation}
\theta_{t+1}=(1-\eta_t\lambda)\theta_t
-\eta_t\frac{\widehat m_t}{\sqrt{\widehat v_t}+\varepsilon}.
\label{eq:adamw}
\end{equation}
若把 $\lambda\theta_t$ 加入梯度，正则项也会改变矩估计并被逐坐标缩放，因而与式~\eqref{eq:adamw} 一般不同。普通无动量 SGD 中，$L_2$ 正则化与相同系数的乘法衰减可以对应；不能把这一等价性直接搬到 Adam。

\section{学习率、正则化与泛化}

学习率预热（warmup）在前若干步逐渐增大学习率，旨在减轻初始梯度尺度和优化器状态尚不稳定时的大更新。余弦衰减可以写成 $\eta_t=\eta_{\min}+\tfrac12(\eta_{\max}-\eta_{\min})(1+\cos(\pi u_t))$，其中 $u_t\in[0,1]$ 是预热后已完成的训练比例。按步数、样本数还是词元数推进调度器会影响不同长度数据下的实际训练，应与预算口径保持一致。

\emph{随机失活（dropout）}在训练时用独立 Bernoulli 掩码 $M_j\sim\operatorname{Bernoulli}(1-p)$ 替换激活为 $M_jh_j/(1-p)$，因此其条件期望仍为 $h_j$；推理时关闭该随机性。这个期望等式不意味着经过后续非线性网络的输出也与训练平均输出完全相同。提前停止、权重衰减和数据增强都可影响泛化，但应由独立验证集选择超参数，测试集只用于最终评估。

对可微、梯度为 $L_s$-Lipschitz 的确定性目标，梯度下降在 $0<\eta\le1/L_s$ 时满足
\[
\mathcal L(\theta-\eta\nabla\mathcal L(\theta))
\le\mathcal L(\theta)-\frac\eta2\|\nabla\mathcal L(\theta)\|_2^2.
\]
其依据是平滑性上界 $\mathcal L(\theta+u)\le\mathcal L(\theta)+\nabla\mathcal L(\theta)^\top u+(L_s/2)\|u\|_2^2$；代入 $u=-\eta\nabla\mathcal L(\theta)$ 并使用 $L_s\eta\le1$ 即得结论。随机梯度、自适应更新与非光滑激活需要另外分析。因此，训练损失总体下降是一项监测结果，而不是所有深度网络更新都满足的逐步保证。

\section{PyTorch 中的计算语义}

PyTorch 的张量保存数值、形状、数据类型和设备信息，自动微分追踪需要梯度的运算。一次训练更新的语义顺序是：得到有效批量，前向计算损失，反向累积梯度，完成同步与必要的缩放，裁剪后更新参数，最后为下一次更新清理梯度。实际程序可以把清零放在循环开头，但不能遗漏；否则无意中改变了优化问题。

\texttt{train()} 与 \texttt{eval()} 控制 dropout 等模块的训练行为；关闭梯度记录控制是否构建求导图，两者用途不同。冻结参数只是不计算该参数的梯度；若其输入还要向前面的可训练层传播梯度，就仍然需要保留相关计算图。保存检查点时，应根据恢复目的保存模型、优化器、调度器、随机数状态和数据进度；只有模型权重足以部署，却不足以精确继续训练。

调试时先在小批量上核对张量形状和目标位置，再验证损失是否能降低。对简单可微模块，可用有限差分抽检导数；对含离散选择、随机失活或低精度算子的程序，必须说明比较条件与容差。训练实现的最小可信单元是损失、梯度和更新之间一致的语义，而不只是一个能运行的训练脚本。

